\section{有理分式的微分和导数}

\begin{proposition}{多项式环上的导子的延拓}{}
	每个$K$-导子$D:K[(X_{i})_{i\in I}]\to K[(X_{i})_{i\in I}]$都可以唯一地延拓为$K$-导子$\bar{D}:K((X_{i})_{i\in I})\to K((X_{i})_{i\in I})$.
\end{proposition}

\begin{remark}{}{}
	注意,这种延拓保持导子之间的括积:若$D_{1},D_{2}$是$K[(X_{i})_{i\in I}]$上的$K$-导子,$\left[D_{1},D_{2}\right]=0$,则$\left[\bar{D}_{1},\bar{D}_{2}\right]=0$.
\end{remark}

\begin{definition}{有理分式的偏导数}{}
	设$f\in K(((X_{i})_{i\in I})$.对每个$i\in I$,称$\mathrm{D}_{X_{i}}f\coloneq\frac{\partial f}{\partial X_{i}}\coloneq f_{X_{i}}^{\prime}\coloneq\bar{\mathrm{D}}_{i}\left(f\right)$是$f$关于不定元$X_{i}$的偏导数.
\end{definition}

\begin{definition}{一元有理分式的导数}{}
	设$f\in K\left(X\right)$.我们称$f$关于$X$的偏导数为$f$关于$X$的导数,记作$\mathrm{D}f$(或$\frac{\mathrm{d}f}{\mathrm{d}X}$,$f^{\prime}$).
\end{definition}

\begin{proposition}{有理分式域的$K$-微分模}{}
	设$B=K[(X_{i})_{i\in I}],C=K((X_{i})_{i\in I})$.考虑典范同构$\Omega_{K}\left(B\right)\otimes_{B}C\to\Omega_{K}\left(C\right)$,则$\Omega_{K}\left(C\right)$有一个基$\left(\mathrm{d}X_{i}\right)_{i\in I}$,并且
	\begin{align*}
		\forall u\in C\left(\mathrm{d}u=\sum\limits_{i\in I}\frac{\partial u}{\partial X_{i}}\mathrm{d}X_{i}\right).
	\end{align*}
	进一步,若$I$是有限集,则$\left(\bar{\mathrm{D}}_{i}\right)_{i\in I}$是$C$-向量空间$\Der_{K}\left(C\right)$的基.
\end{proposition}

\begin{proposition}{导子保持有理分式的代入}{}
	设$E$是$K$-含幺交换结合代数,$\mathbf{x}\in E^{I},f\in K((X_{i})_{i\in I})$.若$\mathbf{x}\in T_{f},y=f\left(\mathbf{x}\right)$,则
	\begin{enumerate}
		\item 对每个满足$\Delta\left(K[(X_{i})_{i\in I}]\right)\subseteq K[(X_{i})_{i\in I}]$的导子$\Delta:K(((X_{i})_{i\in I})\to K(((X_{i})_{i\in I})$有$\mathbf{x}\in T_{\Delta\left(f\right)}$.
		\item 对每个$E$-模$M$和导子$D:E\to M$有$D\left(y\right)=\sum\limits_{i\in I}\frac{\partial f}{\partial X_{i}}\left(\mathbf{x}\right)D\left(x_{i}\right)$.
	\end{enumerate}
\end{proposition}


















































































